\begin{lemma}[Tree-changing transitivity of corresponding ETT presentations]
Fix a finite loopless multigraph \(G\), a palette of \(k\) colors, a
distinguished edge \(e_0\), an input coloring \(\iota\), three exact ordered
displayed trees \(T_0,T_1,T_2\), partial colorings
\(\varphi,\sigma,\tau\), and an extended Tashkinov tree \(\mathcal E_0\)
constructed from \(\iota\) on \(T_0\) and generated by \(\varphi\).  Suppose
\[
 \mathcal C_0:
 \noderef{CorrespondingExtendedTashkinovTree}
   (\mathcal E_0,\sigma,T_1)
\]
is a corresponding presentation, and suppose
\[
 \mathcal C_1:
 \noderef{CorrespondingExtendedTashkinovTree}
   (\mathcal C_0.\mathit{replacementETT},\tau,T_2)
\]
is a corresponding presentation relative to the intermediate replacement
ETT.  Then there is a corresponding presentation
\[
 \mathcal C_{10}:
 \noderef{CorrespondingExtendedTashkinovTree}
   (\mathcal E_0,\tau,T_2)
\]
relative to the original ETT.  Its replacement input is equal to that of
\(\mathcal C_1\), and its dependent replacement series and replacement ETT
are heterogeneously equal to those of \(\mathcal C_1\).
\end{lemma}
\begin{proof}
Write
\[
 n=\mathcal E_0.\mathit{rounds},\qquad
 \mathcal E_m=\mathcal C_0.\mathit{replacementETT},\qquad
 m=\mathcal E_m.\mathit{rounds}.
\]
Let \(\mathcal S_0\) be the replacement series in \(\mathcal C_0\), and
let \(I_1,\mathcal S_1,\mathcal E_1\) be respectively the replacement
input, replacement series, and replacement ETT in \(\mathcal C_1\).
The rung-agreement field of
\noderef{CorrespondingExtendedTashkinovTree} gives an equality
\[
                         h_0:m=n.                              \tag{1}
\]

Transport the round parameter of the dependent
\noderef{TashkinovSeries} \(\mathcal S_1\) along (1), and call the resulting
series \(\widetilde{\mathcal S}_1\).  Define the data fields of
\(\mathcal C_{10}\) to be \(I_1,\widetilde{\mathcal S}_1,\mathcal E_1\).
This choice already gives equality of the two replacement inputs,
heterogeneous equality
\(\widetilde{\mathcal S}_1\mathrel{\mathrm{HEq}}\mathcal S_1\), and
heterogeneous equality of the two copies of \(\mathcal E_1\).  It remains
to construct all the certificate fields.  No equality among
\(T_0,T_1,T_2\) is used: those displayed trees occur only in the types of
the original and replacement ETTs, while every correspondence field below
compares their shared ladder data.

The rung equality for \(\mathcal C_{10}\) is the equality for
\(\mathcal C_1\), followed by (1).  Its dependent series-presentation
field is the series-presentation field of \(\mathcal C_1\), followed by
the canonical heterogeneous equality introduced by transporting
\(\mathcal S_1\) along (1).  Likewise, transporting the final-coloring
field of \(\mathcal C_1\) along (1) proves that the final coloring of
\(\widetilde{\mathcal S}_1\) is exactly \(\tau\).  These transports also
identify the last indices on both sides, rather than identifying merely
their numerical values.

We next verify every indexed structural field.  Fix \(i<n\), and let
\(i_m<m\) be the dependent cast of \(i\) along the inverse of (1).  At
the corresponding positive index, the required ladder-tree equality is
the following explicit chain:
\[
\begin{split}
 \widetilde{\mathcal S}_1.\mathit{tree}(i+1)
 &=\mathcal S_1.\mathit{tree}(i_m+1)\\
 &=\mathcal E_m.\mathit{series}.\mathit{tree}(i_m+1)\\
 &=\mathcal S_0.\mathit{tree}(i+1)\\
 &=\mathcal E_0.\mathit{series}.\mathit{tree}(i+1).
\end{split}                                                    \tag{2}
\]
The first equality is the definition of the dependent transport, the
second is the ladder-tree field of \(\mathcal C_1\), the third is the
series-presentation field of \(\mathcal C_0\) with the same dependent
index cast, and the fourth is the ladder-tree field of \(\mathcal C_0\).
All terms in (2) are exact ordered tree-sequences, so this chain preserves
the individual identities and order of parallel edges.

Now fix any \(a\leq n\), and cast it along (1) to \(a_m\leq m\).  Repeating
the four-term chain (2), first with the appropriate field of
\(\mathcal C_1\), then the series presentation of \(\mathcal C_0\), and
finally the corresponding field of \(\mathcal C_0\), gives
\[
\begin{array}{rcl}
 \widetilde{\mathcal S}_1.\mathit{stored}(a)
  &=&\mathcal E_0.\mathit{series}.\mathit{stored}(a),\\
 \widetilde{\mathcal S}_1.\mathit{connectingEdges}(a)
  &=&\mathcal E_0.\mathit{series}.\mathit{connectingEdges}(a),\\
 \widetilde{\mathcal S}_1.\mathit{kind}(a)
  &=&\mathcal E_0.\mathit{series}.\mathit{kind}(a).
\end{array}                                                     \tag{3}
\]
Thus (3) supplies every raw stored-set, exact connecting-edge-set, and
extension-kind field through the terminal stored index.  The properness
and exactly-one-uncolored fields of \(\mathcal C_1\), transported at the
same indices, supply those two fields for every coloring of
\(\widetilde{\mathcal S}_1\); here properness and unique uncoloredness have
their literal meanings from \noderef{IsProperEdgeColoring} and
\noderef{IsOneEdgeUncolored}.

It remains to prove the positive-stage stability fields.  Fix \(i<n\) and
retain the dependent index \(i_m<m\) used in (2).  The stability field of
\(\mathcal C_0\) says that the \(i+1\)-st coloring of \(\mathcal S_0\) is
\[
 (\mathcal E_0.\mathit{series}.\mathit{tree}(i+1),
  D_i,\mathcal E_0.\mathit{series}.\mathit{coloring}(i+1))
 \text{-stable},                                               \tag{4}
\]
where
\[
 D_i=\noderef{TashkinovStoredColors}
       (\mathcal E_0.\mathit{series}.\mathit{data},i+1).
\]
After transporting through the series-presentation field of
\(\mathcal C_0\), the stability field of \(\mathcal C_1\) says that the
\(i_m+1\)-st coloring of \(\mathcal S_1\) is stable relative to the
\(i+1\)-st coloring of \(\mathcal S_0\), but initially with the tree and
computed protected set belonging to \(\mathcal E_m.\mathit{series}\).

We check, rather than assume, that this latter protected set is \(D_i\).
By \noderef{TashkinovStoredColors}, a positive protected set is the union
of all raw stored sets through the current index, with the colors in the
current tree-missing set deleted.  The raw unions agree by the stored-set
chain in (3).  The tree in question agrees with the original ladder tree
by (2).  Finally, the second clause of
\noderef{StableColoring} in (4) says, vertex by vertex on that exact tree,
that the missing-color set of the \(\mathcal S_0\) coloring equals that of
the original series coloring.  Expanding
\noderef{TreeMissingColors} as the union of these vertexwise missing sets
therefore identifies the two deleted tree-missing sets.  The raw unions
and deleted sets are both equal, so the two computed protected sets are
equal.  Consequently the transported second relation and (4) use the
same exact tree and the same set \(D_i\).

To apply the transitivity clause of \noderef{sc2}, we must place all three
colorings in its subtype of proper colorings leaving exactly \(e_0\)
uncolored.  The middle coloring has these properties by the corresponding
fields of \(\mathcal C_0\), and the final coloring has them by the
corresponding fields of \(\mathcal C_1\).  For the original \(i+1\)-st
coloring, use the certificate at that positive index supplied by
\noderef{TashkinovSeries}.  The initial constructor of
\noderef{TashkinovSeriesStep} cannot occur because the index is positive.
In the SE constructor, its current-coloring equality identifies the
current coloring with the selected coloring, whose properness and
one-uncoloredness are fields of \noderef{MaximumDefectiveChoice}.  In the
PE constructor, the required facts are the constructor's current-proper
and current-one-uncolored fields.  In the RE constructor, its
current-coloring equality identifies the current coloring with the
preceding coloring, and the constructor supplies previous-proper and
previous-one-uncolored for that coloring.  Hence all three subtype
hypotheses hold.

The transitivity field of \noderef{sc2} now composes (4) with the
transported stability relation from \(\mathcal C_1\).  The result is
exactly the \(i\)-th positive replacement-stability field required for
\(\mathcal C_{10}\).  Since \(i<n\) was arbitrary, all such fields have
been constructed.  Together with the rung, series, final-coloring,
structural, properness, and one-uncoloredness fields above, they form the
claimed \(\mathcal C_{10}\), and its three asserted output equalities are
the defining transports chosen in the first paragraph.
\end{proof}
